\chapter{ผลการดำเนินงานด้านงบประมาณ}

\section{ภาพรวมการใช้จ่ายงบประมาณ}
งบประมาณที่ได้รับจัดสรรทั้งสิ้น 7,986,583~บาท โอนถึงผู้รับผิดชอบโครงการแล้ว 7,798,583~บาท
คิดเป็นร้อยละ~98 แสดงว่ากลไกการอนุมัติและโอนงบประมาณของแผนงานทำงานได้ดี
ส่วนที่เป็นคอขวดอยู่ที่ขั้นตอนการใช้จ่ายและเบิกจ่ายในระดับโครงการ

\begin{table}[H]
\centering
\caption{สถานะการใช้จ่ายงบประมาณภาพรวม ณ 4 สิงหาคม 2569}
\small
\begin{tabular}{L{7.4cm}rr}
\toprule
\rowcolor{headerblue}
\textbf{รายการ} & \textbf{จำนวนเงิน (บาท)} & \textbf{ร้อยละ}\\
\midrule
งบประมาณที่ได้รับจัดสรร & 7,986,583 & 100\%\\
\rowcolor{rowblue}
โอนงบประมาณให้ผู้รับผิดชอบแล้ว & 7,798,583 & 98\%\\
ผูกพันใบสั่งซื้อ (ยังไม่เบิกจ่าย) & 825,156 & 10\%\\
\rowcolor{rowblue}
\textbf{เบิกจ่ายจริง} & \textbf{4,460,332} & \textbf{56\%}\\
เบิกจ่ายรวมภาระผูกพัน & 5,285,488 & 66\%\\
\textbf{คงเหลือที่ต้องเบิกจ่าย} & \textbf{3,526,251} & \textbf{44\%}\\
\bottomrule
\end{tabular}
\end{table}

\begin{warnbox}
\textbf{เทียบกับกรอบเวลา} --- ณ 4 สิงหาคม 2569 ระยะเวลาปีงบประมาณผ่านไปแล้ว \textbf{85\%}
ขณะที่เบิกจ่ายได้ \textbf{56\%} ต่างกัน \textbf{29~จุดร้อยละ}
เหลือเวลาถึงวันสิ้นปีงบประมาณ 30 กันยายน 2569 ประมาณ \textbf{55~วัน}
\end{warnbox}

\section{การเบิกจ่ายรายโครงการหลัก}
โครงการหลัก ง8-2 มีอัตราการเบิกจ่ายสูงสุดที่ร้อยละ~64 ขณะที่ ง8-1 และ ง8-3
อยู่ในระดับใกล้เคียงกันคือร้อยละ~47 และ 45 ตามลำดับ โดย ง8-3 เป็นโครงการที่มีวงเงินคงเหลือสูงสุด
เนื่องจากเป็นโครงการที่ใหญ่ที่สุดของแผนงาน

\begin{table}[H]
\centering
\caption{การโอน การผูกพัน และการเบิกจ่ายรายโครงการหลัก (บาท)}
\small
\begin{tabular}{L{3.6cm}rrrrr}
\toprule
\rowcolor{headerblue}
\textbf{โครงการหลัก} & \textbf{จัดสรร} & \textbf{โอนแล้ว} & \textbf{ใบสั่งซื้อ} & \textbf{เบิกจ่าย} & \textbf{คงเหลือ}\\
\midrule
ง8-1 ผลักดันเทคโนโลยี & 1,988,000 & 1,950,000 & 173,100 & 1,064,116 & 669,892\\
\rowcolor{rowblue}
ง8-2 ขับเคลื่อนองค์ความรู้ & 1,999,010 & 1,948,510 & 114,588 & 1,360,534 & 518,312\\
ง8-3 พัฒนากำลังคน & 3,999,573 & 3,900,073 & 443,260 & 2,035,682 & 1,527,517\\
\midrule
\rowcolor{headerblue}
\textbf{รวม} & \textbf{7,986,583} & \textbf{7,798,583} & \textbf{825,156} & \textbf{4,460,332} & \textbf{2,715,721}\\
\bottomrule
\end{tabular}
\end{table}

\begin{infobox}
\textbf{ฐานการคำนวณร้อยละ} --- ร้อยละการเบิกจ่ายทุกจุดในรายงานฉบับนี้คำนวณจาก
\textbf{งบประมาณที่ได้รับจัดสรร} (มิใช่จากยอดที่โอนแล้ว) เพื่อให้ตัวเลขรายโครงการหลัก
บวกกลับได้ตรงกับภาพรวมของแผนงานที่ร้อยละ~50
\end{infobox}

\begin{infobox}
ยอดคงเหลือรายโครงการหลักในตารางข้างต้นคำนวณจาก\textbf{ยอดโอนงบประมาณ}
จึงต่างจากยอดคงเหลือรวมของแผนงาน (3,526,251~บาท) ซึ่งคำนวณจาก\textbf{ยอดงบจัดสรร}
ผลต่าง 188,000~บาท คือส่วนที่ยังไม่ได้โอน
\end{infobox}

\section{การเบิกจ่ายรายหน่วยงาน}
อัตราการเบิกจ่ายระหว่างหน่วยงานมีความแตกต่างกันมาก ตั้งแต่ร้อยละ~93 จนถึงร้อยละ~0
ซึ่งเป็นจุดที่ต้องเข้าไปสนับสนุนเป็นรายหน่วยงาน

\begin{table}[H]
\centering
\caption{งบประมาณและอัตราการเบิกจ่ายรายหน่วยงาน}
\begin{center}\small
\setlength{\tabcolsep}{6pt}\renewcommand{\arraystretch}{1.3}
\begin{tabular}{p{6.6cm}rrrr}
\toprule
\rowcolor{rpfnavy!12}
\textbf{หน่วยงาน} & \textbf{โครงการ} & \textbf{งบจัดสรร} & \textbf{เบิกจ่าย} & \textbf{ร้อยละ}\\
\midrule
กลุ่มแผนงานใต้ร่มพระบารมี & 16 & 2,879,963 & 1,582,954 & 55\%\\
สถาบันถ่ายทอดเทคโนโลยีสู่ชุมชน & 13 & 1,423,650 & 760,822 & 53\%\\
มทร.ล้านนา เชียงราย & 7 & 1,118,960 & 543,709 & 49\%\\
คณะวิศวกรรมศาสตร์ & 6 & 688,000 & 343,314 & 50\%\\
มทร.ล้านนา น่าน & 4 & 555,450 & 532,250 & 96\%\\
วิทยาลัยเทคโนโลยีและสหวิทยาการ & 2 & 349,060 & 275,149 & 79\%\\
มทร.ล้านนา พิษณุโลก & 5 & 331,500 & 265,238 & 80\%\\
มทร.ล้านนา ตาก & 3 & 250,000 & 0 & 0\%\\
คณะศิลปกรรมและสถาปัตยกรรมศาสตร์ & 2 & 200,000 & 96,644 & 48\%\\
สถาบันวิจัยเทคโนโลยีเกษตร & 2 & 150,000 & 60,252 & 40\%\\
คณะบริหารธุรกิจและศิลปศาสตร์ & 1 & 40,000 & 0 & 0\%\\
\midrule
\rowcolor{rpfgold!15}
\textbf{รวม} & \textbf{61} & \textbf{7,986,583} & \textbf{4,460,332} & \textbf{56\%}\\
\bottomrule
\end{tabular}
\end{center}

\end{table}

\begin{warnbox}
\textbf{หน่วยงานที่ยังไม่มีการเบิกจ่าย} --- มทร.ล้านนา ตาก 3~โครงการ วงเงิน 250,000~บาท
และคณะบริหารธุรกิจและศิลปศาสตร์ 1~โครงการ วงเงิน 40,000~บาท
ควรเข้าไปสนับสนุนกระบวนการจัดซื้อจัดจ้างเป็นการเร่งด่วน
\end{warnbox}
